% =========================================================
% Proof: Supremum of a Set Difference
% Source: volume-iii/analysis/bounding/notes/notes-bounds-arithmetic.tex
% =========================================================

\subsection*{Supremum of a Set Difference}
\label{prf:supremum-set-difference}

\begin{remark*}[Return]
$\leftarrow$ \hyperref[thm:sup-diff]{Back to Theorem in Notes}
\end{remark*}

\begin{theorem*}[Supremum of a Difference]
Let $A,B \subset \mathbb{R}$ be nonempty sets that are bounded above and below.
Define
\[
A-B := \{a-b : a\in A,\; b\in B\}.
\]
Then
\[
\sup(A-B) = \sup A - \inf B.
\]
\end{theorem*}

\begin{proof}
Let
\[
C = -B = \{-b : b\in B\}.
\]

Then
\[
A-B = A + C = \{a+c : a\in A,\; c\in C\}.
\]

By the additivity property of the supremum,
\[
\sup(A+C) = \sup A + \sup C.
\]

Next observe that
\[
\sup C = \sup(-B).
\]

Using the scaling rule for negative constants,
\[
\sup(-B) = -\inf B.
\]

Therefore
\[
\sup(A-B)
=
\sup(A+C)
=
\sup A + \sup C
=
\sup A - \inf B.
\]
\end{proof}

\begin{remark*}[Proof shape]
Two-part argument: exhibit the candidate value as an upper (lower) bound,
then show no smaller (larger) value qualifies using the
$\varepsilon$-characterization of the supremum (infimum).
\end{remark*}

\begin{remark*}[Dependencies]
Definition of supremum (infimum); \hyperref[prop:eps-char]{$\varepsilon$-Characterization, Proposition~\ref*{prop:eps-char}}.
\end{remark*}
